\section{Related Work}
\subsection{Target Defense and Reach--Avoid Games}
Cooperative target defense couples interception with protection: defenders must account for both approaching adversaries and the assets they threaten. Differential-game approaches characterize this interaction through geometry, winning regions, and feedback strategies. Garcia et al. derive state-feedback optimal strategies for active defense~\cite{garcia2019activedefense}. For larger teams, Chen et al. combine pairwise reach--avoid outcomes through matching~\cite{chen2017reachavoid}, while Yan et al. use analytical barriers to determine coalition capabilities and formulate task assignment in convex domains~\cite{yan2020barriers}. Matching-based strategies also address three-dimensional heterogeneous teams, combining capture subgames with repeated allocation~\cite{yan2022matching}.

Perimeter defense emphasizes the team outcome when multiple intruders approach a protected region. Shishika et al. develop a decomposition and assignment strategy that explicitly accommodates cooperative defense~\cite{shishika2020perimeter}. This literature establishes that interception opportunities depend on the joint arrangement of defenders and intruders, making the distribution of team effort central to protection. The analytical methods obtain guarantees under their specified dynamics and geometry. Our setting evaluates a learned maneuver policy using cumulative damage to stationary targets, with the same policy deployed across different vehicle and target populations. These changes alter the interactions that the policy must represent and the allocation of effort implicit in its actions.

\subsection{Task Allocation and Cooperative Learning}
Task allocation provides an explicit way to organize multi-robot effort. Gerkey and Matari\'{c} give a formal taxonomy of allocation problems~\cite{gerkey2004allocation}, and consensus-based decentralized auctions combine local bidding with conflict resolution~\cite{choi2009auction}. Cooperative tasks can require robot coalitions: Deng et al. develop distributed auctions for two-robot coalitions, including time-varying assignments and multiplayer reach--avoid applications~\cite{deng2024coalitions}. Miele et al. integrate distributed allocation with safe movement, connectivity maintenance, and deadlock handling~\cite{miele2025coordination}. These approaches expose assignment and motion coordination as coupled decisions.

Learning offers complementary ways to couple these decisions. Centralized training with decentralized execution supports cooperation through multi-agent actor--critic methods and team-reward policy optimization~\cite{lowe2017maddpg,yu2022mappo}. Coach-player learning additionally addresses changing team composition through individual strategies supplied by a coordinating coach~\cite{liu2021copa}. Cooperative distributional reinforcement learning has been applied to dynamic multi-ship weapon--target assignment with attention and delayed rewards~\cite{peng2025wta}. In automation settings, hierarchical learning combines allocation and planning for uncertain swarm confrontation~\cite{wu2025swarm}, and observer-assisted multi-agent learning addresses pursuit--evasion under unknown uncertainties~\cite{liu2025observer}. Closely related to target protection, Li et al. organize learning into intention recognition, target allocation, and path execution~\cite{li2025targetdefense}. Their hierarchy uses differential-game information during allocation training. LEAF instead learns defender action values through a common entity representation, allowing coordination to emerge through maneuver decisions without a separate assignment output. Its contribution concerns the joint relational encoder that supplies these values; the cooperative value-learning objective remains based on REFIL.

\subsection{Entity Representation and Multi-Stage Relational Encoding}
Variable-size entity representations are established tools for cooperative learning. Deep Sets and Set Transformers encode unordered inputs~\cite{zaheer2017deepsets,lee2019settransformer}, while swarm reinforcement learning uses mean embeddings to represent changing neighborhoods~\cite{huttenrauch2019swarm}. Value decomposition, monotonic mixing, and coordination graphs connect agent utilities with team objectives~\cite{sunehag2018vdn,rashid2018qmix,boehmer2020dcg}. REFIL combines entity attention with randomized entity partitions and an auxiliary factorized value objective~\cite{iqbal2021refil}. UPDeT and TransfQMix use transformers in cooperative policies and mixing networks~\cite{hu2021updet,transfqmix2023}; SPECTra emphasizes permutation-consistent scalable networks~\cite{park2026spectra}, and ALMA addresses composite tasks through hierarchical learning~\cite{iqbal2022alma}. Spatial aggregation and temporal recurrence have also been combined for perimeter defense: GR-MAPPO uses GraphSAGE and a GRU and evaluates cross-scale transfer~\cite{tan2026grmappo}.

Relational reinforcement learning models entity interactions with attention~\cite{zambaldi2019relational}. Universal Transformers reuse a transition across depth~\cite{dehghani2019universal}, Jumping Knowledge networks aggregate intermediate graph representations~\cite{xu2018jk}, and LayerDrop trains varying-depth computation paths~\cite{fan2020layerdrop}. Attention-based recurrence addresses temporal information in multi-agent learning~\cite{phan2023aerial}. Building on these precedents, LEAF treats shared contextual updates, self-conditioned intermediate reading, and cross-round fusion as one receiver-conditioned multi-stage relational encoder. One current-self query reads each intermediate entity set, and learned fusion supplies a residual feature to the host's recurrent utility network. The retained GRU carries information across physical decisions. This organization separates contextual processing at a decision from temporal memory across decisions, while jointly learning the loop and its readout for cooperative protection.
